\documentclass[11pt,a4paper,fontset=fandol]{ctexart}

\usepackage[top=24mm,bottom=23mm,left=27mm,right=27mm,headheight=22pt]{geometry}
\usepackage{microtype}
\usepackage{enumitem}
\usepackage{xcolor}
\usepackage{hyperref}
\usepackage{fancyhdr}

\definecolor{ink}{HTML}{1C2530}
\definecolor{softink}{HTML}{626B75}
\definecolor{accent}{HTML}{285FAE}
\definecolor{rulecolor}{HTML}{C9C3B8}

\hypersetup{
  unicode=true,
  pdftitle={Opening Note / 开篇},
  pdfauthor={Xayah Hina},
  pdfsubject={An introduction to the Writing space at i.xayah.me},
  colorlinks=true,
  linkcolor=accent,
  urlcolor=accent
}

\setlength{\parindent}{0pt}
\setlength{\parskip}{0.62em}
\setlist{leftmargin=1.5em,itemsep=0.4em,topsep=0.5em}
\setcounter{secnumdepth}{0}
\setlength{\headheight}{22pt}

\pagestyle{fancy}
\fancyhf{}
\fancyhead[L]{\footnotesize\color{softink}\textsc{Xayah Hina / Writing}}
\fancyhead[R]{\footnotesize\color{softink}Opening Note / 开篇}
\fancyfoot[C]{\footnotesize\color{softink}\thepage}
\renewcommand{\headrulewidth}{0.4pt}
\renewcommand{\headrule}{\hbox to\headwidth{\color{rulecolor}\leaders\hrule height \headrulewidth\hfill}}

\ctexset{
  section={
    format=\Large\bfseries\color{ink},
    beforeskip=1.4em,
    afterskip=0.55em
  },
  subsection={
    format=\normalsize\bfseries\color{accent},
    beforeskip=0.95em,
    afterskip=0.3em
  }
}

\newcommand{\languagepair}[2]{%
  \subsection{English}
  #1
  \subsection{中文}
  #2
}

\begin{document}

\hypersetup{pageanchor=false}
\begin{titlepage}
  \thispagestyle{empty}

  {\small\bfseries\color{accent}\MakeUppercase{Writing No. 001}\par}

  \vspace*{0.18\textheight}

  {\fontsize{39}{43}\selectfont\bfseries\color{ink}Opening Note\par}
  \vspace{0.45em}
  {\fontsize{25}{31}\selectfont\color{softink}开篇\par}

  \vspace{2.2em}
  {\color{rulecolor}\rule{\textwidth}{0.7pt}\par}
  \vspace{1.4em}

  {\large\color{softink}
    A beginning for a place about graphics, music, technical work,
    and the texture of ordinary days.\par
  }
  \vspace{0.8em}
  {\large\color{softink}
    一个关于图形学、音乐、技术工作与普通日常的写作空间，
    从这里开始。\par
  }

  \vfill

  {\large\bfseries Xayah Hina\par}
  \vspace{0.35em}
  {\color{softink}15 July 2026 / 2026 年 7 月 15 日\par}
  \vspace{0.35em}
  {\small\href{https://i.xayah.me/}{i.xayah.me}\par}
\end{titlepage}
\hypersetup{pageanchor=true}

\section{Why This Space Exists / 为什么有这个空间}

\languagepair{
  Some thoughts arrive in the middle of technical work and disappear before
  they have a place to settle. Others begin with a piece of music, a visual
  detail, or an ordinary moment that seems too small for a formal project and
  too important to leave unrecorded. This site is a home for both kinds.

  My academic page is meant to stay concise: research interests, publications,
  projects, and the work that can be named clearly. This space can move at a
  different pace. It can hold unfinished questions, explanations written after
  the code is done, and observations whose value may only become clear later.
}{
  有些想法出现在技术工作的间隙里，还没有找到落脚的地方就消失了。
  另一些想法来自一段音乐、一个视觉细节，或者某个普通得似乎不值得写、
  却又重要得不应该被忘记的瞬间。这个网站想同时容纳这两类东西。

  我的学术主页会尽量保持简洁，只放研究兴趣、论文、项目，以及那些可以被明确命名的工作。
  这里则可以用另一种速度前进：记录尚未完成的问题，写下代码结束之后才出现的解释，
  也保存那些可能要过很久才显出意义的观察。
}

\section{What I Hope to Write / 我想在这里写什么}

The subjects will change, but they will probably return to a few recurring
threads:

\begin{itemize}
  \item \textbf{Graphics and research.} Notes on simulation, rendering,
    inverse design, tools, failed approaches, and the reasoning that rarely
    fits into a project page.
  \item \textbf{Music and listening.} Records of what I am making, hearing,
    and learning from sound.
  \item \textbf{Everyday life.} Small essays, travel notes, photographs in
    words, and moments worth keeping without forcing them into a larger claim.
\end{itemize}

这些主题会不断变化，但大概总会回到几条熟悉的线索：

\begin{itemize}
  \item \textbf{图形学与研究。} 关于模拟、渲染、逆向设计、工具、失败的方法，
    以及那些很难放进项目页面里的推理过程。
  \item \textbf{音乐与聆听。} 记录我正在创作、听见，以及从声音中学到的东西。
  \item \textbf{日常生活。} 短小的随笔、旅途记录、用文字留下的照片，
    以及不必被包装成宏大结论也值得保存的时刻。
\end{itemize}

\section{On Format / 关于形式}

\languagepair{
  Each entry is written in \LaTeX{} and published as a PDF. That choice is
  intentionally a little slower than posting into a feed. Typesetting asks me
  to decide where a thought begins and ends, while the PDF gives each entry a
  stable form that can be read away from the site.

  The web page remains an index rather than becoming another publishing
  system. It tells you what is here; the documents carry the writing itself.
}{
  每篇文章都会用 \LaTeX{} 写成，并以 PDF 发布。这个选择有意比在信息流里发一条内容更慢一些。
  排版迫使我认真决定一个想法从哪里开始、在哪里结束；PDF 也让每篇文章拥有稳定的形态，
  即使离开网站仍然可以被阅读和保存。

  网页本身会保持为一个索引，而不是变成另一套复杂的发布系统。
  它告诉你这里有什么，真正的文字则留在文档中。
}

\section{A Beginning / 一个开始}

\languagepair{
  This first note is deliberately simple. It marks a beginning and leaves
  enough empty space for the character of the site to emerge through future
  writing. If something here resonates with you, you can find my academic work
  at \href{https://graphics.xayah.me/}{graphics.xayah.me}.
}{
  这篇开场有意保持简单。它只是标记一个起点，也为之后的文字留出足够的空白，
  让这个网站的性格慢慢自然形成。如果这里的某些内容与你产生了共鸣，
  也可以在 \href{https://graphics.xayah.me/}{graphics.xayah.me} 找到我的学术工作。
}

\vfill
{\color{rulecolor}\rule{\textwidth}{0.5pt}\par}
{\small\color{softink}Typeset with Tectonic. Source and future notes live at
\href{https://i.xayah.me/}{i.xayah.me}.\par}

\end{document}
